% The raw A_3 (N=4) sweep for the word-level weave
% Delta Delta -> Delta.  This is not the Fock--Goncharov flip: the latter is
% obtained only after applying the cyclic rotation of the boundary chart.
\begingroup
\newcommand{\AThreeSweepCoordinates}[2]{%
  \coordinate (#1top)    at ($(#2)+(0,4)$);
  \coordinate (#1right)  at ($(#2)+(4,0)$);
  \coordinate (#1bottom) at ($(#2)+(0,-4)$);
  \coordinate (#1left)   at ($(#2)+(-4,0)$);
  % Boundary lattice positions.
  \coordinate (#1m13pos)  at ($(#2)+(-1,3)$);
  \coordinate (#1p13pos)  at ($(#2)+(1,3)$);
  \coordinate (#1p22pos)  at ($(#2)+(2,2)$);
  \coordinate (#1p31pos)  at ($(#2)+(3,1)$);
  \coordinate (#1p3m1pos) at ($(#2)+(3,-1)$);
  \coordinate (#1p2m2pos) at ($(#2)+(2,-2)$);
  \coordinate (#1p1m3pos) at ($(#2)+(1,-3)$);
  \coordinate (#1m1m3pos) at ($(#2)+(-1,-3)$);
  \coordinate (#1m2m2pos) at ($(#2)+(-2,-2)$);
  \coordinate (#1m3m1pos) at ($(#2)+(-3,-1)$);
  \coordinate (#1m31pos)   at ($(#2)+(-3,1)$);
  \coordinate (#1m22pos)   at ($(#2)+(-2,2)$);
  % Interior lattice positions.
  \coordinate (#1z2pos)   at ($(#2)+(0,2)$);
  \coordinate (#1m11pos)  at ($(#2)+(-1,1)$);
  \coordinate (#1p11pos)  at ($(#2)+(1,1)$);
  \coordinate (#1m20pos)  at ($(#2)+(-2,0)$);
  \coordinate (#1z0pos)   at ($(#2)+(0,0)$);
  \coordinate (#1p20pos)  at ($(#2)+(2,0)$);
  \coordinate (#1m1m1pos) at ($(#2)+(-1,-1)$);
  \coordinate (#1p1m1pos) at ($(#2)+(1,-1)$);
  \coordinate (#1zm2pos)  at ($(#2)+(0,-2)$);
}

\newcommand{\AThreeSweepTriangleGrid}[3]{%
  \foreach \k in {0,...,4}{%
    \pgfmathsetmacro{\aThreeSweepFraction}{\k/4}
    \draw[aThreeSweepGrid]
      ($(#1)!\aThreeSweepFraction!(#2)$)--
      ($(#1)!\aThreeSweepFraction!(#3)$);
    \draw[aThreeSweepGrid]
      ($(#2)!\aThreeSweepFraction!(#1)$)--
      ($(#2)!\aThreeSweepFraction!(#3)$);
    \draw[aThreeSweepGrid]
      ($(#3)!\aThreeSweepFraction!(#1)$)--
      ($(#3)!\aThreeSweepFraction!(#2)$);
  }
}

% As in the other quiver figures, these are named TikZ nodes rather than
% anonymous marks on coordinates.  TikZ therefore terminates every arrow at
% the node boundary, so no arrowhead is hidden beneath a vertex.
\newcommand{\AThreeSweepSourceNodes}{%
  \foreach \v in {m13,p13,p22,p31,m31,m22}{%
    \node[aThreeSweepFrozen] (L\v) at (L\v pos) {};}
  \foreach \v in {z2,m11,p11,m20,z0,p20,m1m1,p1m1,zm2}{%
    \node[aThreeSweepMutable] (L\v) at (L\v pos) {};}
}

\newcommand{\AThreeSweepTargetNodes}{%
  % The endpoint strings are unchanged by mutation and relabeling: the frozen
  % vertices therefore remain on the two sides 23 and 21 of the upper face.
  \foreach \v in {m13,p13,p22,p31,m31,m22}{%
    \node[aThreeSweepFrozen] (R\v) at (R\v pos) {};}
  \foreach \v in {z2,m11,p11,m20,z0,p20,m1m1,p1m1,zm2}{%
    \node[aThreeSweepMutable] (R\v) at (R\v pos) {};}
}

\newcommand{\AThreeSweepSourceArrows}{%
  \foreach \u/\v in {
    m20/m31,m1m1/m20,m22/m20,m20/m11,zm2/m1m1,
    m11/m1m1,m1m1/z0,p1m1/zm2,p20/p1m1,z0/p1m1,
    p1m1/p11,p31/p20,p11/p20,p20/p22,m11/m22,
    z0/m11,m13/m11,m11/z2,p11/z0,p22/p11,z2/p11,
    p11/p13,z2/m13,p13/z2}{%
    \draw[aThreeSweepArrow] (L\u)--(L\v);}
  \foreach \u/\v in {m31/m22,p22/p31,m22/m13,p13/p22}{%
    \draw[aThreeSweepHalfArrow] (L\u)--(L\v);}
}

\newcommand{\AThreeSweepTargetArrows}{%
  \foreach \u/\v in {
    zm2/m1m1,p1m1/zm2,p1m1/z0,p20/p1m1,p20/p31,
    m1m1/p1m1,m1m1/m20,z0/m1m1,z0/p11,m20/m11,
    m31/m20,m11/z0,m11/m31,m11/z2,p11/p20,
    p11/m11,p11/p22,p31/p11,m22/m11,z2/p11,
    z2/m22,z2/p13,p22/z2,m13/z2,p13/m13}{%
    \draw[aThreeSweepArrow] (R\u)--(R\v);}
  \foreach \u/\v in {m31/m22,m22/m13,p22/p31,p13/p22}{%
    \draw[aThreeSweepHalfArrow] (R\u)--(R\v);}
}

\begin{tikzpicture}[
  x=0.52cm,
  y=0.52cm,
  aThreeSweepGrid/.style={draw=black!16,line width=0.34pt},
  aThreeSweepBoundary/.style={draw=black!56,line width=0.72pt,
    line cap=round,line join=round},
  aThreeSweepDiagonal/.style={draw=wqorange,line width=0.86pt,
    line cap=round},
  aThreeSweepArrow/.style={wqlocalarrow,line width=0.50pt,
    -{Stealth[length=1.18mm,width=0.84mm]}},
  aThreeSweepHalfArrow/.style={aThreeSweepArrow,wqhalfedge},
  aThreeSweepMutable/.style={circle,wqqvertex,fill=white,
    minimum size=1.90mm,inner sep=0pt,line width=0.52pt},
  aThreeSweepFrozen/.style={rectangle,wqqvertex,fill=teal!6,
    minimum size=1.90mm,inner sep=0pt,line width=0.52pt},
  aThreeSweepPanel/.style={font=\scriptsize,text=black!82},
  aThreeSweepMap/.style={wqmaparrow},
  aThreeSweepSequence/.style={font=\scriptsize,text=black!88,
    align=left,inner sep=0pt}
]
\path[use as bounding box] (-13.8,-5.35) rectangle (13.8,4.8);
\coordinate (leftCenter) at (-9,0);
\coordinate (rightCenter) at (9,0);
\AThreeSweepCoordinates{L}{leftCenter}
\AThreeSweepCoordinates{R}{rightCenter}

% Initial boundary: the quadrilateral is divided by the vertical diagonal.
\AThreeSweepTriangleGrid{Lleft}{Lbottom}{Ltop}
\AThreeSweepTriangleGrid{Lbottom}{Lright}{Ltop}
% Final boundary: the quadrilateral is divided by the horizontal diagonal.
\AThreeSweepTriangleGrid{Rleft}{Rright}{Rtop}
\AThreeSweepTriangleGrid{Rleft}{Rbottom}{Rright}
\draw[aThreeSweepBoundary]
  (Ltop)--(Lright)--(Lbottom)--(Lleft)--cycle;
\draw[aThreeSweepBoundary]
  (Rtop)--(Rright)--(Rbottom)--(Rleft)--cycle;
\draw[aThreeSweepDiagonal] (Ltop)--(Lbottom);
\draw[aThreeSweepDiagonal] (Rleft)--(Rright);

% Put down named nodes first, as in the manuscript's other quiver figures;
% arrows then meet the node boundaries automatically.
\AThreeSweepSourceNodes
\AThreeSweepTargetNodes
\AThreeSweepSourceArrows
\AThreeSweepTargetArrows

\node[aThreeSweepPanel,anchor=north] at ($(leftCenter)+(0,-4.55)$)
  {initial $Q(\boldsymbol\Delta\boldsymbol\Delta)$};
\node[aThreeSweepPanel,anchor=north] at ($(rightCenter)+(0,-4.55)$)
  {final $Q^{\mathrm{sweep}}(\fW)$};

\draw[aThreeSweepMap] (-3.35,0)--(3.35,0);
\node[font=\small,fill=white,inner xsep=1.5pt,inner ysep=0.8pt]
  at (0,0.58) {sweep $\mu_\fW$};
\node[font=\scriptsize,fill=white,inner xsep=1.2pt,inner ysep=0.6pt]
  at (0,-0.57)
  {$\boldsymbol\Delta\boldsymbol\Delta\longrightarrow\boldsymbol\Delta$};

% A top-to-bottom sweep of the tetrahedral drawing, reading equal-height
% vertices from left to right.  Four-valent commutations give no mutation.
% \node[aThreeSweepSequence,anchor=north] at (0,-5.55) {$
% \begin{aligned}
%  \mu_\fW:\quad
%  &\mu_{(0,0,2,0)}\ ;\ \mu_{(0,0,1,1)}\ ;\
%   \mu_{(1,0,1,0)}\ ;\ \mu_{(0,1,1,0)}\ ;\\[-0.10ex]
%  &\mu_{(0,0,0,2)}\ ;\ \mu_{(1,0,0,1)}\ ;\
%   \mu_{(0,1,0,1)}\ ;\\[-0.10ex]
%  &\mu_{(2,0,0,0)}\ ;\ \mu_{(1,1,0,0)}\ ;\
%   \mu_{(0,2,0,0)}.
% \end{aligned}$};
\end{tikzpicture}
\endgroup
